% ============================================================
%  Artifact Description (AD) Appendix -- REVISED DRAFT
%  SC 2026
%
%  Revision rationale: the original appendix described a
%  limited-access *API* to the MIST engine, justified by
%  "proprietary intellectual property that cannot be shared
%  publicly".  MIST is now fully open source, so that carve-out
%  is removed and the artifact becomes two conventional
%  open-source repositories.  Reproduction times below are
%  MEASURED, not estimated.
% ============================================================

\documentclass[conference]{IEEEtran}
\IEEEoverridecommandlockouts

\usepackage{amsmath,amssymb,amsfonts}
\usepackage{graphicx}
\usepackage{hyperref}
\usepackage{booktabs}
\usepackage{multirow}
\usepackage{array}
\usepackage{tabularx}
\usepackage{xcolor}
\usepackage{enumitem}
\usepackage{geometry}
\geometry{margin=0.8in, top=0.8in}

\newcommand{\tool}{{MIST }}
\newcommand{\toolns}{{MIST}}

\begin{document}

\title{Artifact Description Appendix: MIST --- A Co-Design Framework for
Heterogeneous, Multi-Stage LLM Inference}

\author{\normalsize{Supercomputing 2026}}

\maketitle

% ============================================================
\section*{Part 1: Overview of Contributions \& Artifacts}
% ============================================================

\subsection*{A. Paper's Main Contributions}

The paper makes the following contributions for which computational
artifacts are provided:

\begin{enumerate}[leftmargin=1.2em]

  \item \textbf{End-to-end LLM inference simulation framework (\tool).}
        \tool is an event-driven simulator that models the full serving
        stack --- workload, system/software, and hardware layers --- for
        heterogeneous, multi-stage LLM inference pipelines.

  \item \textbf{Heterogeneous hardware design-space exploration (DSE).}
        Using \toolns, we systematically search multi-vendor prefill-decode
        (PD) disaggregated deployment configurations and demonstrate that
        heterogeneous configurations achieve higher generated tokens per
        second per dollar than homogeneous single-vendor baselines.

  \item \textbf{KV cache storage co-design.}
        \tool enables hardware co-design for KV cache storage by
        evaluating five distinct memory-hierarchy architectures,
        revealing how storage granularity and bandwidth interact with
        end-to-end tail latency under varying KV reuse patterns.

\end{enumerate}

\subsection*{B. Computational Artifacts}

\begin{table}[h]
\centering
\caption{Computational Artifacts Overview}
\label{tab:artifacts}
\resizebox{\columnwidth}{!}{%
\begin{tabular}{@{}p{2.6cm} p{3.6cm} p{1.8cm} p{3.2cm}@{}}
\toprule
\textbf{Artifact} & \textbf{Description} & \textbf{Related Contribution} & \textbf{Reproduced Paper Elements} \\
\midrule
\tool Simulator (A1) &
  The complete, open-source \tool simulation engine.
  Accepts request traces, model specifications, and hardware
  parameters; returns per-request and aggregate latency,
  throughput, and cost metrics. &
  C1, C2, C3 &
  All experiments \\
\addlinespace
Reproduction Scripts (A2) &
  Python scripts that drive \toolns, sweep the relevant
  parameter spaces, collect raw results, and generate
  publication-ready figures. &
  C1, C2, C3 &
  Fig.~4 (a--d), Fig.~6 (a--b), Fig.~7 (a--b), Fig.~8, Fig.~10 \\
\bottomrule
\end{tabular}%
}
\end{table}

\noindent \textbf{A1 --- \tool simulator:}
\begin{center}
\url{https://github.com/MIST-Simulator/MIST}
\end{center}

\noindent \textbf{A2 --- Reproduction scripts:}
\begin{center}
\url{https://github.com/MIST-Simulator/SC_Paper_Charts}
\end{center}

\noindent Both are released under the MIT license.  A2 declares A1 as a
pinned dependency in \texttt{requirements.txt}, so installing A2 installs
the exact \tool revision used to produce the paper's figures.

% ============================================================
\section*{Part 2: Artifact Identification}
% ============================================================

\subsection*{A. Relation to Contributions}

The primary artifact is the \tool simulator itself, released in full.
It accepts:
\begin{itemize}[leftmargin=1.2em]
  \item \textbf{Workload input:} real or synthetic request traces
        (arrival times, prompt/decode token lengths, and optional
        prefix-KV context lengths).
  \item \textbf{Model specification:} model name and tensor-parallelism
        degree.
  \item \textbf{Hardware configuration:} accelerator SKU, cluster
        topology, and memory-hierarchy parameters.
\end{itemize}
\tool returns per-request and aggregate metrics as described in
\S III.G, including Time to First Token (TTFT), Time per Output Token
(TPOT), Time to Last Token (TTLT), step-level runtimes, and simulation
wall-clock time.

The reproduction scripts (A2) drive \tool across all experimental
configurations described in the paper, write raw per-configuration
result files, and invoke a separate set of plotting scripts to
regenerate the figures.

\paragraph{Reproducibility of simulated data.}
Every \toolns-derived data point in every figure is produced by running
the simulator from the scripts in A2; no simulated value is
hand-transcribed.  Request traces are generated or replayed under fixed
random seeds and \toolns's hardware models are deterministic, so
re-running a task reproduces its result CSV, and therefore its figure,
exactly.

\paragraph{Vendored baseline data.}
Measurements that cannot be regenerated without the original hardware or
third-party simulators ship as data files in A2, each documented with its
provenance: measured vLLM runs (L40S$\times$2, H100$\times$8, TPUv6e),
Vidur outputs, LLMServingSim~2.0 outputs, and \texttt{fio} block-device
latencies.  These are the reference points \tool is validated
\emph{against}, not \toolns outputs.

\subsection*{B. Expected Results}

\paragraph{Validation (Fig.~4, Fig.~6).}
\tool reproduces measured on-hardware end-to-end latency distributions
with a mean absolute error of 6.5--11.5\% across the three validation
platforms, while simulating 40--290$\times$ faster than the baseline
simulators (Fig.~4(d)).
For the per-step benchmarks (Fig.~6(a)), under a held-out evaluation
split \toolns's ML-based LLM cluster model attains 2.1\% (prefill),
0.5\% (decode), and 2.1\% (chunked prefill) mean absolute error, versus
18.3\%, 7.7\%, and 14.9\% for Vidur.
For KV retrieval (Fig.~6(b)), \tool tracks \texttt{fio} measurements to
7.7\% (NVMe SSD) and 14.4\% (DDR4) for block sizes at or above 16\,MB;
relative error is larger below 4\,MB, where absolute latencies are small.
These results substantiate Contribution~1.

\paragraph{Heterogeneous DSE (Fig.~7, Fig.~8).}
The search confirms that heterogeneous multi-vendor PD-disaggregated
configurations Pareto-dominate homogeneous single-vendor alternatives on
the generated tokens/\$ axis, for both the chat and the code-generation
use cases.
Absolute tokens/\$ depends on the cloud rental prices listed in \S5.1
(sampled November 2025); the performance ranking depends only on the
hardware models and is unaffected by price movement.
These results substantiate Contribution~2.

\paragraph{KV Cache Storage DSE (Fig.~10).}
The five storage architectures exhibit clearly differentiated latency
CDFs.  Dedicated per-client memory (Case~A) and platform-level shared
memory (Case~B) deliver the lowest tail latency for private KV caches at
both context lengths, while rack-level shared SSD (Case~C) is best for
shared global caches with long KV sequences.
These results substantiate Contribution~3.

\subsection*{C. Reproduction Time}

Times below are measured on a 2023-class laptop-grade workstation
(10 CPU cores, 32\,GB RAM), not estimated.

\begin{table}[!th]
\centering
\caption{Measured Reproduction Times}
\label{tab:times}
\begin{tabular}{@{}lp{4.6cm}@{}}
\toprule
\textbf{Phase} & \textbf{Time} \\
\midrule
Artifact Setup      & $\sim$15 min (environment + trace download) \\
Figure regeneration & $<$1 min (from shipped reference results) \\
\addlinespace
T1 execution        & $\sim$1 min \\
T2 execution        & $\sim$1 min \\
T4 execution        & $\sim$22 min \\
T3 execution        & 6--8 h (\texttt{--fast-eval}); 20+ h (full) \\
\midrule
\textbf{Total}      & \textbf{$\sim$7--9 h (reduced T3); 21+ h (full)} \\
\bottomrule
\end{tabular}
\end{table}

\noindent A reviewer who only wishes to confirm that the figures follow
from the recorded results can regenerate all seven figures in under a
minute with \texttt{bash run\_all.sh --plot-only}; the shipped reference
results are themselves the output of the \texttt{run\_T*.py} invocations
recorded in \texttt{results/reference/PROVENANCE.md}.

\subsection*{D. Artifact Setup}

\subsubsection*{Hardware}

A standard Linux or macOS workstation with 8--16 CPU cores and at least
16\,GB of RAM.  No GPU, accelerator, or special interconnect is required:
\tool is an analytical/event-driven simulator and runs entirely on CPU.
An internet connection is needed to install dependencies and download
the input traces.

\subsubsection*{Software}

\begin{table}[h]
\centering
\caption{Required Software}
\label{tab:software}
\resizebox{\columnwidth}{!}{%
\begin{tabular}{@{}llll@{}}
\toprule
\textbf{Package} & \textbf{Version} & \textbf{Purpose} & \textbf{URL} \\
\midrule
Conda / Miniconda & $\geq$23.x &
  Environment management &
  \url{https://docs.conda.io} \\
Python            & 3.10--3.12 &
  Experiment and plotting scripts &
  \url{https://www.python.org} \\
\tool             & pinned in \texttt{requirements.txt} &
  Simulation engine &
  See A1 above \\
GenZ              & pinned in \texttt{requirements.txt} &
  Analytical hardware model used by \tool &
  \url{https://github.com/abhibambhaniya/GenZ-LLM-Analyzer} \\
\bottomrule
\end{tabular}%
}
\end{table}

\noindent All remaining Python dependencies (\texttt{numpy},
\texttt{pandas}, \texttt{scipy}, \texttt{scikit-learn},
\texttt{matplotlib}, \texttt{seaborn}) are pinned in
\texttt{requirements.txt} and installed by the setup script.

\subsubsection*{Datasets / Input}

\begin{enumerate}[leftmargin=1.2em]

  \item \textbf{Microsoft Azure LLM Inference Traces (2023).}
        Public request traces used by the KV cache storage study.
        {\small\url{https://github.com/Azure/AzurePublicDataset/blob/master/AzureLLMInferenceDataset2023.md}}

  \item \textbf{ShareGPT conversation dataset.}
        Used by the end-to-end validation experiments (T1).

\end{enumerate}

\noindent \texttt{scripts/download\_traces.sh} fetches both and places
them in the expected directory structure.  The derived per-use-case
traces required by T3 are committed to A2 and need no download.  The
measured baseline data (vLLM, Vidur, LLMServingSim~2.0, \texttt{fio})
is likewise committed to A2.

\subsubsection*{Installation \& Deployment}

\begin{enumerate}[leftmargin=1.2em]
  \item Clone A2:
        \texttt{git clone https://github.com/MIST-Simulator/SC\_Paper\_Charts.git}
  \item Run \texttt{./setup.sh} to create the Conda environment and
        install all dependencies, including \tool itself.
  \item Activate it: \texttt{conda activate mist-sc26}.
  \item Download the traces:
        \texttt{bash scripts/download\_traces.sh}.
  \item Verify the installation:
        \texttt{python scripts/smoke\_test.py}
        (imports \toolns, runs a short simulation, and renders a figure;
        completes in under a minute).
\end{enumerate}

\subsection*{E. Artifact Evaluation}

The workflow consists of four independent task groups, which may be run
sequentially or concurrently (\texttt{bash run\_all.sh --parallel}).

\medskip
\noindent\textbf{T1 --- End-to-End Runtime Validation}
\hfill\textit{($\sim$1 min | Reproduces: Fig.~4 a--d)}

\begin{itemize}[leftmargin=1.2em, topsep=2pt]
  \item \textbf{T1.1} --- Qwen3-32B on L40S$\times$2 (TP2).
        Reproduces Fig.~4(a) and Fig.~4(d).
  \item \textbf{T1.2} --- Llama3-70B on H100$\times$8 (TP8).
        Reproduces Fig.~4(b) and Fig.~4(d).
  \item \textbf{T1.3} --- Llama-3.1-8B on TPUv6e.
        Reproduces Fig.~4(c) and Fig.~4(d).
\end{itemize}
Each replays the same fixed 100-request ShareGPT trace, at the same
arrival schedule, that the vLLM/Vidur/LLMServingSim baselines were
measured against; holding the workload fixed is what makes the
comparison controlled.  \texttt{--resample-sharegpt} draws a fresh
seeded sample instead, for sensitivity analysis.\\
\textit{Dependencies:} none.
\smallskip

\noindent\textbf{T2 --- Individual Step Validation}
\hfill\textit{($\sim$1 min | Reproduces: Fig.~6 a--b)}

\begin{itemize}[leftmargin=1.2em, topsep=2pt]
  \item \textbf{T2.1} --- Per-step prefill/decode/chunked runtimes for
        Llama-2-70B on H100 (TP4 and TP8), evaluated on a held-out
        split.  Reproduces Fig.~6(a).
  \item \textbf{T2.2} --- KV cache retrieval latency for NVMe SSD and
        DDR4 across block sizes from 256\,KB to 1\,GB, against
        \texttt{fio} measurements.  Reproduces Fig.~6(b).
\end{itemize}
\textit{Dependencies:} none.
\smallskip

\noindent\textbf{T3 --- Heterogeneous Design-Space Search}\\
\textit{(6--8 h reduced; 20+ h full | Reproduces: Fig.~7 a--b, Fig.~8)}

\begin{itemize}[leftmargin=1.2em, topsep=2pt]
  \item \textbf{T3.1} --- DSE for the chat conversation use case
        (no prefix-KV reuse), at 100\,RPS with a TTFT P99 SLO.
        Reproduces Fig.~7(a) and the chat column of Fig.~8.
  \item \textbf{T3.2} --- DSE for the code generation use case with high
        prefix-KV reuse, at 100\,RPS with a relaxed TTFT P99 SLO.
        Reproduces Fig.~7(b) and the code-generation column of Fig.~8.
\end{itemize}
Both search hardware SKU $\times$ tensor parallelism $\times$ batching
strategy $\times$ prefill-to-decode client ratio, and rank
SLO-satisfying configurations by generated tokens/second/\$.\\
\textit{Reduced search-space option:} \texttt{--fast-eval} restricts the
search to a deterministic, stratified $\sim$150-configuration subset per
use case that spans the Pareto frontier.  The sweep is resumable:
re-running after an interruption skips completed configurations.\\
\textit{Dependencies:} none; T3.1 and T3.2 are independent.
\smallskip

\noindent\textbf{T4 --- KV Cache Storage Design-Space Search}
\hfill\textit{($\sim$22 min | Reproduces: Fig.~10)}

Evaluates the five KV cache storage architectures of Table~2 on a
256-GPU system (128 clients, H100:TP2, 4 racks), sweeping short (4K
token) and long (24K token) KV contexts for both private and shared
cache scenarios.\\
\textit{Dependencies:} none.

\subsection*{F. Artifact Analysis}

Each task group follows a two-script pipeline:

\begin{enumerate}[leftmargin=1.2em]
  \item \textbf{Experiment script}
        (\texttt{run\_\{T1$|$T2$|$T3$|$T4\}.py}):
        drives \tool over the configured parameter sweep, collects raw
        per-request and aggregate metrics, and writes structured
        \texttt{.csv} files to \texttt{results/<task>/}.
  \item \textbf{Plotting script}
        (\texttt{plot\_\{T1$|$T2$|$T3$|$T4\}.py}):
        reads those CSVs, applies post-processing (CDFs, normalised
        tokens/\$, overlaid real-hardware reference points), and writes
        \texttt{.pdf} figures to \texttt{figures/}.
        The plotting scripts do not import \toolns, so figures can be
        regenerated without the simulator installed.
\end{enumerate}

\noindent Task dependency graph:

\begin{center}
  T1.1, T1.2, T1.3 $\;\rightarrow\;$ \texttt{plot\_T1.py}
    $\;\rightarrow\;$ Fig.~4 \\[3pt]
  T2.1, T2.2 $\;\rightarrow\;$ \texttt{plot\_T2.py}
    $\;\rightarrow\;$ Fig.~6 \\[3pt]
  T3.1, T3.2 $\;\rightarrow\;$ \texttt{plot\_T3.py}
    $\;\rightarrow\;$ Fig.~7, Fig.~8 \\[3pt]
  T4 $\;\rightarrow\;$ \texttt{plot\_T4.py}
    $\;\rightarrow\;$ Fig.~10
\end{center}

\noindent \texttt{run\_all.sh} executes all four experiment scripts
(sequentially, or concurrently with \texttt{--parallel}) and then all
four plotting scripts.  \texttt{run\_all.sh --plot-only} regenerates
every figure from existing results without re-simulating.

\end{document}
